\documentclass[11pt]{article}
\usepackage{amsmath,amssymb,amsfonts,amsthm}
\usepackage{booktabs}
\usepackage{hyperref}
\usepackage{geometry}
\geometry{margin=1in}

\newtheorem{theorem}{Theorem}[section]
\newtheorem{lemma}[theorem]{Lemma}
\newtheorem{proposition}[theorem]{Proposition}
\newtheorem{definition}[theorem]{Definition}
\newtheorem{remark}[theorem]{Remark}
\newtheorem{corollary}[theorem]{Corollary}

\newcommand{\R}{\mathbb{R}}
\newcommand{\Z}{\mathbb{Z}}

\title{Uniqueness-First Effective Field Theory Realization\\
of Logarithmic Mass Lattices:\\
Framework, Feasibility, and the Spectrum Obstruction}
\author{Marvin L.\ Gentry}
\date{}

\begin{document}
\maketitle

\begin{abstract}
We present a minimal effective field theory framework in which fermion mass
hierarchies are organized by a discrete logarithmic lattice. Each mass
arises from a unique leading operator, fixed by a shaping symmetry and a
finite set of spurions and their conjugates. We conduct a rigorous
feasibility analysis of the hypothesis that the nine charged Standard Model
fermion masses lie on a common three-generator integer lattice with
well-conditioned basis and small integer coordinates. The analysis identifies
a concrete obstruction: the near-coincidence of the strange-quark and muon
masses ($|\ln m_s-\ln m_\mu|\approx0.124$ natural units) combined with the
large dynamic range of the full spectrum ($\approx12.8$ natural log units)
forces either large coordinates, badly conditioned generators, or trivially
decoupled subsectors. Within the charged-lepton sector alone, a
one-generator fit with integer coordinates $(0,2,3)$ achieves a residual
of $5.8\%$ on the $\tau$ mass. We report these as the best results
currently achievable and identify what additional structure would be
required to extend them.
\end{abstract}

\section{Introduction}

Fermion masses in the Standard Model span roughly thirteen orders of
magnitude. In logarithmic coordinates, nontrivial regularities motivate a
careful organizational analysis. We ask whether these regularities are
captured by a discrete integer lattice: a finite set of basis vectors
$\{\lambda_i\}$ such that each fermion's logarithmic mass is well
approximated by $\sum_i n_i\lambda_i$ with small integers $n_i$.

This paper proceeds in two parts. Sections~\ref{sec:eft}--\ref{sec:cp}
develop the EFT framework independently of any specific numerical lattice.
Section~\ref{sec:feasibility} tests the simplest lattice hypothesis against
PDG data and proves a precise obstruction theorem.

\section{Uniqueness-First Operator Realization}
\label{sec:eft}

\subsection{Setup}

We work in logarithmic coordinates $q=\alpha\ln(m/m_0)$. We hypothesize
basis vectors $\{\lambda_1,\lambda_2,\lambda_3\}\subset\R_{>0}$ and integer
assignments $(a_f,b_f,c_f)\in\Z^3$ such that
\begin{equation}\label{eq:lattice}
q(m_f)\approx a_f\lambda_1+b_f\lambda_2+c_f\lambda_3.
\end{equation}

\subsection{Spurion framework}

A shaping symmetry $\mathcal{S}$ forbids renormalizable mass terms. Scalar
spurions $\{\Phi_i\}$ carry charges $Q_i$ under $\mathcal{S}$. Fermion
masses arise from operators
\begin{equation}\label{eq:operator}
\mathcal{O}_f=\frac{1}{\Lambda^{d-4}}\bar\psi_L H\psi_R
\prod_i\Phi_i^{\,n_i}(\Phi_i^\dagger)^{\bar n_i},
\end{equation}
giving $m_f\sim v\prod_i(\langle\Phi_i\rangle/\Lambda)^{n_i-\bar n_i}$.
Negative coordinates ($n_i-\bar n_i<0$) are realized by net conjugate
spurion insertions without inverse operators.

\subsection{Uniqueness}

\begin{theorem}[Uniqueness of leading operator]\label{thm:uniqueness}
Given fixed $\mathcal{S}$ and spurion set, the leading operator for each
fermion $f$ is unique.
\end{theorem}

\begin{proof}
The leading operator minimizes total spurion count $\sum_i(n_i+\bar n_i)$
subject to charge neutrality. Any other contribution involves additional
$\Phi_i\Phi_i^\dagger$ pairs and is suppressed by additional powers of
$\langle\Phi_i\rangle/\Lambda<1$.
\end{proof}

\section{Sector-by-Sector Analysis}
\label{sec:sectors}

Table~\ref{tab:coords} records the integer coordinate assignments. The
charged-lepton entries are established by the single-generator fit of
Section~\ref{sec:feasibility}. Quark entries are undetermined pending
resolution of the spectrum obstruction (Theorem~\ref{thm:obstruction}).

\begin{table}[h]
\centering
\begin{tabular}{llccc}
\toprule
Sector & Fermion & $a_f$ & $b_f$ & $c_f$ \\
\midrule
Charged leptons & $e$ & 0 & 0 & 0 \\
 & $\mu$ & 2 & 0 & 0 \\
 & $\tau$ & 3 & 0 & 0 \\
\midrule
Up-type quarks & $u$ & $a_u$ & $b_u$ & $c_u$ \\
 & $c$ & $a_c$ & $b_c$ & $c_c$ \\
 & $t$ & $a_t$ & $b_t$ & $c_t$ \\
\midrule
Down-type quarks & $d$ & $a_d$ & $b_d$ & $c_d$ \\
 & $s$ & $a_s$ & $b_s$ & $c_s$ \\
 & $b$ & $a_b$ & $b_b$ & $c_b$ \\
\bottomrule
\end{tabular}
\caption{Integer lattice coordinates. Lepton entries from
Proposition~\ref{prop:leptons}; quark entries undetermined
(see Theorem~\ref{thm:obstruction}).}
\label{tab:coords}
\end{table}

\begin{remark}[Neutrino sector]\label{rem:neutrinos}
Neutrino masses lie at $\ln(m_\nu/m_e)\approx-16$ to $-20$, requiring
large negative coordinates or a separate mechanism. We leave the neutrino
sector to a separate treatment.
\end{remark}

\section{Discrete CP/T Structure}
\label{sec:cp}

The plaquette invariant
\begin{equation}\label{eq:plaquette}
\Theta=\arg\!\left(\frac{Y_{12}Y_{23}Y_{31}}{Y_{13}Y_{21}Y_{32}}\right)
\end{equation}
is invariant under all independent rephasings. Since Yukawa phases belong
to a finite alphabet determined by $\mathcal{S}$, $\Theta$ is restricted
to a finite set: CP/T violation is discrete.

The \emph{cap inversion} map $\phi_{ij}\mapsto-\phi_{ij}$ sends
$\Theta\mapsto-\Theta$. By Lemma~3.1 of the companion paper on CP phases
from hyperbolic holonomy, orientation reversal of a closed geodesic induces
exactly this sign flip on holonomy phases. Cap inversion is therefore the
algebraic counterpart of geometric orientation reversal.

\section{Feasibility Analysis: The Spectrum Obstruction}
\label{sec:feasibility}

\subsection{Data}

\begin{table}[h]
\centering
\begin{tabular}{llrr}
\toprule
Sector & Fermion & $\ln(m_f/m_e)$ & $\ln(m_f/m_e)/\ln\varphi$ \\
\midrule
Charged leptons & $e$ & 0.000 & 0.000 \\
 & $\mu$ & 5.333 & 11.082 \\
 & $\tau$ & 7.944 & 16.507 \\
\midrule
Up quarks & $u$ & 1.441 & 2.994 \\
 & $c$ & 7.813 & 16.235 \\
 & $t$ & 12.818 & 26.637 \\
\midrule
Down quarks & $d$ & 2.212 & 4.596 \\
 & $s$ & 5.209 & 10.824 \\
 & $b$ & 8.999 & 18.698 \\
\bottomrule
\end{tabular}
\caption{Logarithmic mass coordinates relative to $m_e=0.51100$~MeV.
PDG~2022~\cite{PDGreview} central values.}
\label{tab:masses}
\end{table}

\subsection{The obstruction theorem}

\begin{theorem}[Spectrum obstruction]\label{thm:obstruction}
Let $\delta_{\min}=0.124$ (the $s$--$\mu$ separation) and
$q_{\max}=12.818$ (the top quark). Any three-generator integer lattice fit
to all nine charged fermion masses with residual $\varepsilon<\delta_{\min}/2$
and coordinate bound $N$ satisfies
\begin{equation}\label{eq:condbound}
\kappa:=\frac{\max_i\lambda_i}{\min_i\lambda_i}
\geq\frac{q_{\max}}{3N\,\delta_{\min}}
=\frac{34.46}{N}.
\end{equation}
For $N\leq6$ the basis is ill-conditioned ($\kappa>5$).
\end{theorem}

\begin{proof}
Since $\varepsilon<\delta_{\min}/2$, distinct fermions map to distinct
lattice points. The $s$--$\mu$ pair differs by $0.124$, forcing a nonzero
lattice vector of magnitude $<0.248$, so $\lambda_{\min}\leq0.248$. The
top quark at $q_{\max}=12.818$ with coordinates bounded by $N$ forces
$\lambda_{\max}\geq12.818/(3N)$. Dividing gives~\eqref{eq:condbound}.
\end{proof}

\subsection{Single-sector result}

\begin{proposition}[Charged-lepton fit]\label{prop:leptons}
With $\lambda_1=\ln(m_\mu/m_e)/2=2.6665$ and coordinates $(0,2,3)$:
the $e$ and $\mu$ are exact; the $\tau$ residual is
$|3\times2.6665-7.9440|=0.0555$, a mass ratio error of $5.7\%$.
The midpoint fit $\lambda_1=2.657$ gives $\pm1\%$ on each.
\end{proposition}

\section{Discussion and Outlook}

The uniqueness-first EFT framework and the spectrum obstruction theorem are
the primary results. The obstruction rules out the simplest discrete
generalization of Froggatt--Nielsen~\cite{FroggattNielsen} across all three
sectors simultaneously. The most structurally motivated path forward is
sector-specific generator subsets drawn from a common eight-dimensional
shape-space basis (see companion paper on discrete shape space), with
cross-sector near-coincidences explained by algebraic relations between
generators of different sectors. Connecting cap inversion explicitly to the
holonomy structure of the CP companion paper is a natural next step.

\section{Conclusion}

We have established the uniqueness-first EFT framework, the spectrum
obstruction theorem, and the best currently achievable single-sector
lattice result. The framework is falsifiable and we have reported the
results of the first feasibility test honestly. Extending to the full
spectrum requires additional structure identified in the obstruction proof.

\bibliographystyle{plain}
\bibliography{paper7}

\end{document}